\section*{Как использовались ИИ-агенты}
Автор выполнил измерения и передал три страницы журнала, дополнительное
описание и фотографию установки. Codex организовал перенос данных,
расчёт и подготовку LaTeX. Агент --- отдельный рабочий чат модели
с конкретным заданием и доступом к файлам проекта.

\begin{center}
\begin{tikzpicture}[x=1cm,y=1cm,
 block/.style={rounded corners=3pt,text=white,align=center,
 font=\bfseries,text width=3.5cm,minimum height=1.05cm,inner sep=6pt},
 flow/.style={->,line width=1pt,draw=labblue}]
 \node[block,fill=labblue] (raw) at (0,0) {Журнал};
 \node[block,fill=labgreen] (csv) at (4.6,0) {Данные и расчёт};
 \node[block,fill=labred] (pdf) at (9.2,0) {Отчёт};
 \draw[flow] (raw) -- (csv);
 \draw[flow] (csv) -- (pdf);
 \node[draw=labblue,fill=labblue!7,rounded corners=3pt,
 text width=9cm,align=center,inner sep=8pt,font=\small]
 (check) at (4.6,-1.8) {\textbf{Повторное чтение $\;\cdot\;$ Пересчёт $\;\cdot\;$ Проверка PDF}};
 \draw[flow] (raw.south) -- (raw.south |- check.north);
 \draw[flow] (csv.south) -- (check.north);
 \draw[flow] (pdf.south) -- (pdf.south |- check.north);
\end{tikzpicture}
\end{center}

\begin{table}[h!]\centering\small
\caption{Участники и задачи. Идентификаторы моделей проверены по журналам запусков.}
\begin{tabular}{p{2.3cm}p{3cm}p{9cm}}
\hline
Агент & Модель OpenAI & Задача\\ \hline
Ведущий & GPT-6 Astra & Постановка задач, повторный ввод чисел, независимые вычисления, сборка и публикация.\\
\texttt{data} & GPT-6 Astra & Перенос 39 разрядных и 83 зондовых измерений в CSV, знаки по П2, отметки исправлений.\\
\texttt{processing} & GPT-6 Astra & МНК, физические параметры, проверки устойчивости, модуль нового бюджета погрешностей.\\
\texttt{report} & GPT-6 Astra & Методика, таблицы, графики, физическое обсуждение и исходная вёрстка.\\
\texttt{review} & GPT-6 Luna & Проверка формул, источников, выводов и реализации нового бюджета погрешностей.\\ \hline
\end{tabular}
\end{table}

\paragraph{Перенос и проверка рукописи.}
Агент \texttt{data} читал увеличенные страницы и записывал пары напряжения
и тока. Для П2 «минус» обращал знаки обеих координат. Ведущий независимо
ввёл пары из скана и сравнил их с CSV программно. Расхождения разобрали
по увеличенным ячейкам. Десять первоначально сомнительных прочтений
отметили как \texttt{uncertain}. Затем автор проверил и подтвердил все
значения. Метки и опыт временного исключения этих точек сохранены
как история проверки. В итоговой обработке используются все 83 точки.
Фото и уточнение автора подтвердили $d=0{,}2$ мм и $l=5{,}2$ мм.

\paragraph{Как получили результаты.}
Агент \texttt{processing} заменил графические построения п.~9 методички
тремя прямыми МНК. Центральная прямая описывает наклон около нуля,
две крайние --- ток насыщения. Их пересечения определяют сдвиг координат,
ионный ток и энергию электронов. Концентрация и параметры экранирования
вычисляются по формулам пп.~11--14 в СИ. Дифференциальное сопротивление
находится по окнам из трёх соседних разрядных точек.

\clearpage
\section*{Как проверялись вычисления и погрешности}
\paragraph{Независимое воспроизведение.}
Ведущий написал \texttt{verify\_results.py}, который заново строит прямые
через \texttt{numpy.polyfit}, находит их пересечения и рассчитывает параметры
плазмы. Для разрядной ВАХ использована \texttt{scipy.stats.linregress}.
Проверяются все три режима и четыре экстремальных наклона.
Проверяющая программа читает исходные таблицы и сравнивает свои числа
с результатами обработки.

\paragraph{Новый бюджет погрешностей.}
По указанию автора приняты пределы в одну единицу последнего разряда
и 3\,\% показания прибора. Пересчёт поручили тому же агенту
\texttt{processing}; он подготовил \texttt{ai\_uncertainty.py}.
После остановки помощников по лимиту использования ведущий завершил
подключение модуля, проверку, графики и текст этой версии.
Пределы представлены равномерными распределениями. Общая погрешность
шкалы одного прибора действует сразу на все его показания.
Это сохраняет важные сокращения, например множитель шкалы тока
сокращается в отношении ионного тока к центральному наклону.

При распространении ошибок программа учитывает обе координаты каждой точки,
ковариацию параметров МНК и связь измеренных точек с вычисленным центром
характеристики. К приборному бюджету добавлен остаточный разброс регрессий.
Допущения о корреляциях, номиналах установки и характере распределений
приведены в разделе о погрешностях.

\paragraph{Проверка нового модуля.}
Ведущий написал \texttt{verify\_ai\_uncertainty.py} без импорта модуля
погрешностей. Он поочерёдно немного изменяет исходные напряжения и токи,
каждый раз заново строит прямые и вычисляет результат. Из откликов
восстанавливает вклад каждого разряда и каждой общей шкалы.
Так проверены 39 бюджетов производных величин, четыре наклона разрядной ВАХ
и напряжение зажигания. Отдельные тесты проверяют сокращение общей шкалы
микроамперметра в энергии и 27 элементов межрежимных ковариаций
энергии, концентрации и тока разряда по законам масштабирования.
Повторный запуск обработки
проверен сравнением файлов.

\paragraph{Проверка устойчивости и смысла.}
Центральное окно меняли с 4,6 В на 3 и 6,5 В, крайние границы ---
с 16 В на 15 и 18 В. Также давали двум полярностям отдельные свободные
члены при общем наклоне. При 4,99 мА последнее изменение повышает
энергию с 3,49 до 5,20 эВ. Поэтому расширенный бюджет приборных
ошибок сохраняет зависимость результата от модели.
Рецензент первой версии проверил формулы Бома, единицы,
классификацию режима и библиографию по методичке и практикуму.
Рецензент повторно проверил новый модуль и указал на необходимость
явного учёта общей шкалы тока разряда между режимами.
После дополнения матрицы ведущий проверил её отдельной программой.

\paragraph{Публикация и журнал действий.}
PDF просматриваются после сборки, таблицы и графики формируются из CSV.
GitHub Actions повторяет обработку и сборку из опубликованных исходников.
Установленный обработчик журналов вызывается вручную для локальных копий
чатов. Проверяются идентификаторы, непустые JSONL с вызовами инструментов
и читаемые копии Markdown. Это позволяет проследить задания,
изменения файлов и результаты проверок.
